\chapter{Droites parallèles}

\begin{itemize}
\it Démontrer qu'une perpendiculaire $(Ax)$ et une oblique $(By)$ à une même droite $(AB)$ sont sécantes.
\it Démontrer que si deux droites $(\Delta)$ et $(\Delta')$ sont respectivement perpendiculaires à deux droites parallèles $(D)$ et $(D')$, elles sont parallèles entre elles. 
\it Démontrer que, si deux droites $(Ax)$ et $(By)$ sont respectivement perpendiculaires aux côtés d'un angle saillant $\widehat{AOB}$, elles sont concourantes. 
\it Soit un triangle isocèle $ABC$ de base $[BC]$. \begin{enumerate}
\item Montre rque la bissectrice extérieure de l'angle $\widehat{A}$ est parallèle à la base $[BC]$.
\item La réciproque est-elle vraie ? La démontrer.
\end{enumerate}
\it Soient $[AB]$ et $[CD]$ deux diamètres d'un cercle de centre $O$.
\begin{enumerate}
\item Comparer les directions de $[AC]$ et de $[BD]$ avec celle de la bissectrice de l'angle $\widehat{AOC}$.
\item Que peut-on en conclure pour $(AC)$ et $(BD)$ ?
\end{enumerate}
\it Soit un angle $\widehat{xOy}$. Du point $O$ comme centre, on trace un premier cercle qui coupe $[Ox)$ en $A$ et $[Oy)$ en $B$, puis un second cercle qui coupe $[Ox)$ en $C$ et $[Oy)$ en $D$. 
\begin{enumerate}
\item Comparer les directions de $[AB]$ et de $[CD]$ avec celle de la bissectrice de l'angle $\widehat{xOy}$.
\item En déduire que $(AB)$ et $(CD)$ sont parallèles.
\end{enumerate}
\it On considère sur un cercle de centre $O$ qutre points $A$, $B$, $C$ et $D$ dans cet ordre et tels que les arcs $\arc{AB}$ et $\arc{CD}$ soient égaux. 
\begin{enumerate}
\item Montrer que les angles $\widehat{AOD}$ et $\widehat{BOC}$ ont même bissectrice.
\item Comparer les directions de $(AD)$ et de $(BC)$.
\end{enumerate}
\it D'un point intérieur $M$ on mène les perpendiculaires $(MA)$ et $(MB)$ aux côtés de l'angle
droit $\widehat{xOy}$.
\begin{enumerate}
\item Montre que $(MA)$ et $(Oy)$ sont parallèles, et qu'il en est de même de $(MB)$ et $(Ox)$.
\item Quelle est la valeur de l'angle $\widehat{AMB}$ ?
\end{enumerate}
\it On considère un point $M$ pris sur un demi-cercle de diamètre $[AB]$ et de centre $O$. On mène les perpendiculaires $(OH)$ à $(MA)$ et $(OK)$ à $(MB)$.
\begin{enumerate}
\item Montrer que l'angle $\widehat{HOK}$ est droit. Comparer les directions de $(OH)$ et de $(MB)$ ainsi que celles de $(OK)$ et de $(MA)$.
\item Évaluer l'angle $\widehat{AMB}$ et montrer que les angles $\widehat{MAB}$ et $\widehat{MBA}$ sont complémentaires.
\end{enumerate}
\it Soit un quadrilatère $ABCD$ dans lequel les diagonales $[AC]$ et $[BD]$ ont même milieu $O$. On mène les perpendiculaires $(OH)$ à $(AB)$ et $(OK)$ à
$(CD)$. \begin{enumerate}
\item Comparer les triangles $OAB$ et $OCD$, puis les triangles $OAH$ et $OCK$. Conséquences ?
\item Comment sont disposés les points $O$, $H$ et $K$ ? En déduire que $(AB)$ et $(CD)$ sont parallèles. En est-il de même pour $(AD)$ et $(BC)$ ?
\end{enumerate}
\it Étant donné une droite $(xy)$ et un point extérieur $O$, on mène la perpendiculaire $(OH)$ et une oblique $(OM)$ à la droite $(xy)$ puis on construit $H'$ et $M'$ symétriques de $H$ et $M$ par rapport à $O$. \begin{enumerate}
\item Comparer les triangles $OHM$ et $OH'M'$. Conséquences ?
\item Que peut-on dire des directions $(HM)$ et $(H'M')$ ? En déduire que : \\
\emph{Deux droites symétriques l'une de l'autre par rapport à un point sont parallèles.}
\item En déduire une construction de la parallèle menée par $A$ à la droite $(xy)$.
\end{enumerate}
\it On considère deux parallèles $(D)$ et $(D')$ perpendiculaires en $A$ et $A'$ au segment $[AA']$ dont on désigne par $O$ le milieu.\begin{enumerate}
\item Démontrer que la médiatrice $(xy)$ du segment $[AA']$ est un axe de symétrie de la bande $(D, D')$.
\item Établir que $(xy)$ est médiatrice de tout segment $[HH']$ perpendiculaire à $(D)$ et $(D')$.
\item En déduire que tout point $I$ de $(xy)$ est un centre de symétrie des deux droites $(D)$ et $(D')$ et qu'il est équidistant de ces deux droites.
\end{enumerate}
\end{itemize}